\section{Related Works}

Recent advances in Large Language Models (LLMs) and agent-based systems have demonstrated significant potential in the field of Autonomous Scientific Research (ASR), enabling progress from creative idea generation to end-to-end research automation. Some studies~\citep{li2024chain, wang2023scimon, zhou2024hypothesis} have shown that LLMs are capable of generating novel research ideas, which has sparked widespread discussion in the academic community. For example, \citet{li2024chain} introduce a method that derives research ideas through the analysis of interconnected scholarly works. Beyond idea generation, several studies have examined the use of LLMs for hypothesis formulation~\citep{qi2023large, qi2024large}, such as extracting hypotheses from large-scale web data~\citep{yang2023large}, multi-agent framework using LLMs to enhance collaborative hypothesis generation in biomedicine~\citep{qi2023large, qi2024large}, and scientific literature~\citep{wang2023scimon, zhou2024hypothesis}. However, most of these efforts remain at the stage of idea or hypothesis generation, lacking systematic empirical validation of their practical effectiveness.

In terms of end-to-end research automation, ~\citet{lu2024ai} introduced the AI Scientist framework, which was among the first to achieve a fully automated pipeline in the machine learning domain, covering problem definition, experimental execution, and result reporting. The subsequent AI Scientist-V2~\citep{yamada2025ai} further enhanced the framework by incorporating agent tree search, vision-language model feedback, and parallelized experiment execution, leading to the first workshop paper fully generated and peer-reviewed by AI. Similarly, systems such as AI-Researcher~\citep{AiResearcher} and Dolphin~\citep{yuan2025dolphin} have proposed closed-loop, LLM-driven frameworks that automate the entire research process on a range of simple tasks.

Human-AI collaboration is gaining traction in ASR. Systems like Agent Laboratory~\citep{Agentlaboratory} integrate human feedback into multi-stage LLM agent workflows, automating literature review, experiment execution, and report writing, while allowing user input at each step to enhance research quality. AgentRxiv~\citep{Agentrxiv} addresses the collaborative nature of scientific discovery by enabling LLM agent laboratories to communicate and build upon each other's work via a shared preprint server, thus facilitating knowledge sharing and collective innovation. Experimental results demonstrate that agent laboratories utilizing AgentRxiv for collaboration achieve greater performance improvements compared to isolated settings. Similarly, AI Co-Scientist~\citep{gottweis2025towards}, based on Gemini 2.0, employs a multi-agent system with a "generate-debate-evolve" strategy for hypothesis generation, and has demonstrated effectiveness in biomedical domains such as drug repurposing, novel target identification, and interpretation of bacterial evolution, with several hypotheses validated through experiments.

Despite these advances, most current systems are still evaluated primarily on relatively simple tasks or within narrow scientific domains. However, when applied to more complex, system-level scientific challenges, these approaches often face significant limitations. Key challenges include generating truly novel and scientifically sound research ideas, establishing robust closed-loop feedback between experiments and idea generation, and developing systematic evaluation standards to rigorously assess the effectiveness and real-world value of autonomous research systems.